%==============================================================
%  PHẦN ĐẦU: Lời cam đoan, Lời cảm ơn, Danh mục từ viết tắt
%==============================================================

\chapter*{Lời cam đoan}
\addcontentsline{toc}{chapter}{Lời cam đoan}
\markboth{Lời cam đoan}{Lời cam đoan}

Tôi xin cam đoan báo cáo này do tôi thực hiện, dưới sự hướng dẫn của
\thesisAdvisor. Các kết quả lý thuyết được trình bày ở đây hoặc do tôi tự chứng
minh, hoặc được trích dẫn đầy đủ từ tài liệu gốc; những chỗ tôi chỉ diễn giải
lại một kết quả đã có đều được ghi rõ là như vậy. Toàn bộ số liệu thực nghiệm
được tạo ra bằng mã nguồn công bố kèm báo cáo (Phụ lục~\ref{app:code}) và có thể
kiểm chứng lại; các giới hạn về tính tái lập của những số liệu ấy được nêu thẳng
ở Mục~\ref{sec:hanche}. Báo cáo không sao chép từ bất kỳ công trình nào khác mà
không ghi rõ nguồn.

\vspace{1.5cm}
\begin{flushright}
\begin{minipage}{8.6cm}
\centering
TP. Hồ Chí Minh, \thesisDate\\[0.3cm]
\textit{Sinh viên thực hiện}\\[2cm]
\thesisAuthor
\end{minipage}
\end{flushright}

%--------------------------------------------------------------
\chapter*{Lời cảm ơn}
\addcontentsline{toc}{chapter}{Lời cảm ơn}
\markboth{Lời cảm ơn}{Lời cảm ơn}

Trước hết, tôi xin bày tỏ lòng biết ơn sâu sắc tới \thesisAdvisor{}, người đã
nhận hướng dẫn và theo sát tôi trong suốt quá trình thực hiện báo cáo này. Thầy
không chỉ gợi ý đề tài mà còn chỉ cho tôi \emph{cách} tiếp cận nó: danh mục tài
liệu tham khảo chi tiết mà thầy giao, với chỉ dẫn đọc đến từng mục của từng công
trình, là thứ đã định hình toàn bộ cấu trúc của báo cáo. Quan trọng hơn, thầy
đã nhiều lần yêu cầu tôi phân biệt rạch ròi giữa điều đã chứng minh và điều mới
chỉ được suy đoán. Yêu cầu ấy khó chịu hơn tôi tưởng lúc đầu, và cuối cùng lại
trở thành mạch chính của cả chương thực nghiệm: phần lớn những gì báo cáo thu
được nằm ở chỗ dám ghi lại các kết quả không đạt kỳ vọng thay vì bỏ qua chúng.

Tôi xin cảm ơn quý Thầy Cô Khoa Toán--Tin học, Trường Đại học Khoa học Tự nhiên,
Đại học Quốc gia Thành phố Hồ Chí Minh. Những kiến thức về giải tích, phương
trình đạo hàm riêng, giải tích số và tối ưu hoá mà tôi được học ở Khoa là nền
tảng trực tiếp cho báo cáo này; nhiều chỗ trong Chương~\ref{chap:pinn} và
Chương~\ref{chap:thucnghiem} chỉ viết được vì tôi đã có sẵn những công cụ ấy.
Tôi cũng cảm ơn Thầy Cô trong hội đồng đã dành thời gian đọc và góp ý cho bản
báo cáo này.

Xin cảm ơn bạn bè cùng khoá đã trao đổi, tranh luận và nhiều lần chỉ ra những
chỗ tôi lập luận vội trong các bản nháp trước.

Cuối cùng, tôi xin gửi lời cảm ơn tới gia đình, những người đã tạo mọi điều kiện
và kiên nhẫn với tôi trong suốt thời gian học tập.

Báo cáo chắc chắn còn thiếu sót. Tôi mong nhận được góp ý của quý Thầy Cô và
các bạn để có thể hoàn thiện hơn.

\vspace{1cm}
\begin{flushright}
\begin{minipage}{8.6cm}
\centering
TP. Hồ Chí Minh, \thesisDate\\[0.3cm]
\thesisAuthor
\end{minipage}
\end{flushright}

%--------------------------------------------------------------
\chapter*{Danh mục từ viết tắt}
\addcontentsline{toc}{chapter}{Danh mục từ viết tắt}
\markboth{Danh mục từ viết tắt}{Danh mục từ viết tắt}

\begin{longtable}{@{}l p{5.3cm} p{6.0cm}@{}}
\toprule
\textbf{Viết tắt} & \textbf{Nguyên văn} & \textbf{Nghĩa dùng trong báo cáo} \\
\midrule
\endfirsthead
\toprule
\textbf{Viết tắt} & \textbf{Nguyên văn} & \textbf{Nghĩa dùng trong báo cáo} \\
\midrule
\endhead
\bottomrule
\endfoot
PINN & physics-informed neural network & mạng nơ-ron thông tin vật lý \\
FNN & feedforward neural network & mạng nơ-ron truyền thẳng \\
MLP & multilayer perceptron & perceptron nhiều lớp \\
JVP & Jacobian--vector product & tích Jacobi--vectơ (chế độ thuận) \\
VJP & vector--Jacobian product & tích vectơ--Jacobi (chế độ ngược) \\
MSE & mean squared error & sai số bình phương trung bình \\
SGD & stochastic gradient descent & hạ gradient ngẫu nhiên \\
BFGS & Broyden--Fletcher--Goldfarb--Shanno & phương pháp tựa Newton \\
L-BFGS & limited-memory BFGS & BFGS bộ nhớ hạn chế \\
ReLU & rectified linear unit & hàm kích hoạt $\max\{0,z\}$ \\
SiLU & sigmoid linear unit & hàm kích hoạt Swish, $z\,\sigma_{\mathrm{s}}(z)$ \\
LHS & Latin hypercube sampling & lấy mẫu siêu khối Latin \\
RAR & residual-based adaptive refinement & lấy mẫu thích nghi theo phần dư \\
NTK & neural tangent kernel & nhân tiếp tuyến nơ-ron \\
RK4 & Runge--Kutta bậc bốn & lược đồ tích phân thời gian tường minh \\
cPINN & conservative PINN & PINN phân rã miền dạng bảo toàn \\
XPINN & extended PINN & PINN phân rã miền tổng quát \\
\bottomrule
\end{longtable}

\begin{center}
\begin{minipage}{0.92\textwidth}
\small
\textbf{Về ba chữ viết tắt không dùng.} Báo cáo viết đầy đủ ``phương trình đạo
hàm riêng'' và ``phương trình vi phân thường'' thay vì PDE và ODE, để tránh lẫn
với cách viết tắt tiếng Việt; và viết ``vi phân tự động'' thay vì AD, dù AD là
lối viết chuẩn trong tài liệu tiếng Anh --- ghi ở đây để người đọc đối chiếu
được. Riêng ký hiệu $\mathcal{L}$ thì luôn chỉ toán tử vi
phân và không bao giờ chỉ hàm mất mát (xem quy ước ở đầu
Chương~\ref{chap:pinn}).
\end{minipage}
\end{center}
